\documentclass[11pt]{article}
\usepackage[a4paper,margin=27mm]{geometry}
\usepackage{amsmath,amssymb,amsthm,xcolor,booktabs,array}
\newcommand{\PROVED}{\textcolor{green!40!black}{\textbf{PROVED}}}
\newcommand{\OPEN}{\textcolor{red!70!black}{\textbf{OPEN}}}
\newcommand{\AUDIT}{\textcolor{orange!70!black}{\textbf{AUDIT}}}
\newtheorem{theorem}{Theorem}
\newtheorem{proposition}{Proposition}
\title{The half-Tate Banach envelope and row (d) (v101)}
\date{13 August 2026}

\begin{document}
\maketitle

\section{Metric linearization of row (b)}

Define
\begin{equation}
 \mathcal B_{1/2}:=left\{a=(a_n)_{n\ge1}:
       \|a\|_{1/2}:=\sum_{n\ge1}|a_n|n^{-1/2}<\infty\right\}             \tag{1}
\end{equation}
with Dirichlet convolution.  This is not chosen from the zeros or the Weil
form: its weight is exactly the self-dual metric character
$w_{1/2}(\Gamma_n)=n^{-1/2}$ of row (b).

\begin{theorem}[Half-Tate Banach algebra --- \PROVED]
The space $\mathcal B_{1/2}$ is a commutative unital Banach algebra and
\begin{equation}
             \|\delta_m*\delta_n\|_{1/2}
             =\|\delta_m\|_{1/2}\|\delta_n\|_{1/2}.             \tag{2}
\end{equation}
The algebraic row-(b) correspondence linearization therefore extends
canonically to (1).
\end{theorem}

\begin{proof}
The weight is completely multiplicative:
$(mn)^{-1/2}=m^{-1/2}n^{-1/2}$.  Hence
\[
 \|a*b\|_{1/2}
 \le\sum_{m,n}|a_m||b_n|(mn)^{-1/2}
 =\|a\|_{1/2}\|b\|_{1/2}.
\]
Completeness is ordinary weighted $\ell^1$ completeness, the unit is
$\delta_1$, and equality (2) is immediate.
\end{proof}

\section{Extension of the individual-correspondence functor}

Row (c) sends $\delta_n$ to $U_n f(x)=f(nx)$ on
$\mathcal S_>=\mathcal S_{(1,\infty)}$.  Its defining estimate is, for
every relevant seminorm at an exponent $a>1$,
\[
 q_{a,j,k}(U_nf)\le Cn^{-a}(1+\log n)^N
        \sum_{j'\le j}q_{a,j',k}(f).                 \tag{3}
\]

\begin{theorem}[Half-Tate multiplier extension --- \PROVED]
For every $a=(a_n)\in\mathcal B_{1/2}$, the series
\begin{equation}
                  \mathcal R_{1/2}(a)f:=\sum_{n\ge1}a_nU_nf             \tag{4}
\end{equation}
converges in $\mathcal S_>$, uniformly on bounded subsets in every defining
multiplier seminorm.  It defines a continuous algebra homomorphism
\begin{equation}
 \mathcal R_{1/2}:\mathcal B_{1/2}longrightarrow
                     \mathsf{Mult}(\mathcal S_>).     \tag{5}
\end{equation}
\end{theorem}

\begin{proof}
For $a>1$ and fixed $N$,
$n^{-a}(1+\log n)^N\le C_{a,N}n^{-1/2}$.  Multiplying (3) by $|a_n|$ and
summing proves absolute convergence controlled by
$C_{a,N}\|a\|_{1/2}$.  Dirichlet convolution and $U_mU_n=U_{mn}$ prove the
homomorphism identity by absolute rearrangement.
\end{proof}

Thus the metric line of row (b) does produce a genuine norm-enriched
extension of the row-(c) functor on its individual-correspondence target.

\section{Character test}

On the simple Mellin character with parameter $s$,
\begin{equation}
                         \chi_s(\delta_n)=n^{-s}.     \tag{6}
\end{equation}

\begin{theorem}[Automatic continuity test --- \PROVED]
The algebraic character (6) extends to a continuous character of
$\mathcal B_{1/2}$ if and only if
\begin{equation}
                              \Re s\ge\frac12.        \tag{7}
\end{equation}
When it extends, its norm is one.
\end{theorem}

\begin{proof}
A character on a unital Banach algebra is continuous and has norm one.
Necessity also follows directly from point masses:
\[
 |n^{-s}|\le\|\delta_n\|_{1/2}=n^{-1/2}
 \quad\Longrightarrow\quad \Re s\ge1/2.
\]
Conversely, under (7),
\[
 \left|\sum_na_nn^{-s}\right|
 \le\sum_n|a_n|n^{-\Re s}
 \le\|a\|_{1/2},
\]
which defines the required continuous character and gives norm one because
it takes the unit to one.
\end{proof}

\begin{corollary}[Banach-envelope criterion for row (d) --- \PROVED]
If the semisimplified odd row-(c) action extends from its algebraic
$\mathbb C[\mathbb N^\times]$ action to a continuous
$\mathcal B_{1/2}$-module action compatible with (5), then every spectral
parameter has real part $1/2$ and row (d) follows.
\end{corollary}

\begin{proof}
Every simple quotient of a continuous unital Banach module is a continuous
character, so (7) gives $\Re\rho\ge1/2$.  The Fourier--Poisson involution of
row (c) also supplies $1-\bar\rho$; applying (7) to it gives
$1-\Re\rho\ge1/2$.  Thus $\Re\rho=1/2$.  The trace-preserving Loewy
assembly of v99 then gives the negative Hilbert square for primitive test
functions.
\end{proof}

\section{Does the extension act on the odd quotient?}

The answer is not automatic.  Although $\mathcal H_-\subset\mathcal S_>$,
the series (4) need not preserve $\mathcal H_-$.  For example, choose a
nonzero smooth $f$ supported near $1$ and coefficients
$a_n=n^{-1/2-\varepsilon}$ with $0<\varepsilon\le1/2$.  Then
$a\in\mathcal B_{1/2}$, but near the points $x=n^{-1}$ the term
$a_n f(nx)$ has size comparable to $n^{-1/2-\varepsilon}$.  Hence the sum
has only polynomial decay as $x\downarrow0$ and is not in
$\mathcal S_{(-\infty,\infty)}=\mathcal H_-$.

Moreover the Zeta element is not in the half-Tate algebra:
\begin{equation}
              \left\|\sum_{n\ge1}\delta_n\right\|_{1/2}
              =\sum_{n\ge1}n^{-1/2}=\infty.          \tag{8}
\end{equation}
Therefore the existence of (5) does not by itself make
$\mathcal H_-^0=\mathcal H_-/Z\mathcal H_\cap$ a
$\mathcal B_{1/2}$-module.

The precise remaining extension theorem is
\begin{equation}
 \boxed{
 \text{the algebraic action on }\operatorname{gr}_L(\mathcal H_-^0)
 \text{ is continuous for }\|\cdot\|_{1/2}.}         \tag{9}
\end{equation}
It is enough to establish (9) on every finite Loewy quotient, uniformly in
the quotient.  By the character test and functional symmetry, (9) closes
row (d) immediately.

Statement (9) is \OPEN.  It is not a new normalization: (1)--(5) construct
the normalization and its multiplier action completely.  It is the
boundary-descent theorem saying that the half-Tate multiplier action does
not lose continuity when passing from the one-sided space $\mathcal S_>$
to the two-sided Poisson cokernel.

\begin{theorem}[Boundary descent is the exact terminal theorem ---
\PROVED]
On the semisimplified odd object, (9) is equivalent to condition $G$ (and
hence to row (d) under the reductions already proved in the manuscript).
It is not stronger than the desired conclusion and it is not a merely
technical continuity lemma.
\end{theorem}

\begin{proof}
The corollary above proves (9)$\Rightarrow$ every spectral parameter is
critical, hence $G$ by the row-(c) spectral formula and the supplied Weil
criterion.  Conversely, under $G$ every simple parameter has
$\Re\rho=1/2$.  The character test shows that each simple action extends
contractively to $\mathcal B_{1/2}$; taking the Loewy semisimplification
and its multiplicity direct sum gives (9).  The nilpotent jets are absent
from the semisimplification but their multiplicities are retained by the
finite trace theorem of v99.
\end{proof}

\begin{center}
\begin{tabular}{>{\raggedright\arraybackslash}p{.57\linewidth}
                >{\centering\arraybackslash}p{.15\linewidth}
                >{\raggedright\arraybackslash}p{.18\linewidth}}
\toprule
Statement & Status & Location \\
\midrule
Half-Tate Banach algebra & \PROVED & (1)--(2) \\
Norm-enriched functor on $\mathcal S_>$ & \PROVED & (3)--(5) \\
Continuous characters exactly satisfy $\Re s\ge1/2$ & \PROVED & (6)--(7) \\
Functional symmetry upgrades this to $\Re s=1/2$ & \PROVED & corollary \\
Descent of the Banach action to the odd Loewy object & \OPEN & (9) \\
Row (d), condition $G$ & \OPEN & immediate after (9) \\
\bottomrule
\end{tabular}
\end{center}

\end{document}
